\documentclass[10pt,a4paper]{article}
\usepackage[margin=1.5cm,top=1.2cm,bottom=1.2cm]{geometry}
\usepackage{multicol}
\usepackage{booktabs}
\usepackage{enumitem}
\usepackage{listings}
\usepackage{xcolor}
\usepackage[T1]{fontenc}
\usepackage{lmodern}
\usepackage{microtype}
\usepackage{titlesec}

\pagestyle{empty}
\setlength{\parindent}{0pt}
\setlength{\parskip}{3pt}
\setlength{\columnsep}{0.6cm}

\titleformat{\section}{\bfseries\normalsize}{}{0pt}{}
\titlespacing*{\section}{0pt}{4pt}{2pt}
\titleformat{\subsection}{\bfseries\small}{}{0pt}{}
\titlespacing*{\subsection}{0pt}{3pt}{1pt}

\lstset{
  basicstyle=\ttfamily\scriptsize,
  backgroundcolor=\color{gray!10},
  frame=single,
  framerule=0.3pt,
  xleftmargin=2pt,
  xrightmargin=2pt,
  aboveskip=3pt,
  belowskip=3pt,
  breaklines=true,
}

\begin{document}

\begin{center}
{\large\bfseries Measuring and Improving Multi-Target Binary Selection\\
in LLVM's GPU Offload Stack}\\[3pt]
{\normalsize S.\ Akash\quad \textit{IIT Patna\enspace$\cdot$\enspace CERN GSoC alumnus\enspace$\cdot$\enspace vLLM contributor}}
\end{center}

\vspace{-2pt}
\hrule height 0.5pt
\vspace{4pt}

\begin{multicols}{2}

\section{Thesis}

MLIR can compile a single \texttt{gpu.module} to N GPU targets (NVIDIA, AMD, Intel),
and \texttt{OffloadBinary} can carry all N images in one fat binary.
But at runtime, LLVM's offload stack picks the \emph{first compatible image} and stops---there
is no standard metadata for ``best-compatible'' selection, and no published measurement of
what this costs.

\section{Contributions}

\begin{enumerate}[leftmargin=1.2em,itemsep=1pt,topsep=1pt]
\item \textbf{Standard metadata vocabulary} for \texttt{OffloadBinary}:
  5 new keys enabling runtime selection beyond triple/arch matching
  (additive, backward-compatible, $\sim$30 LOC header patch).
\item \textbf{First published flame graph} of the LLVM GPU dispatch stack:
  per-layer latency from \texttt{OffloadBinary} parse through
  \texttt{cuLaunchKernel} on a GTX~1650 (10{,}000 iterations, 100 cold trials).
\item \textbf{Design sketch} for \texttt{\#gpu.runtime\_select}:
  an MLIR attribute deferring binary selection from compile time to runtime.
\end{enumerate}

\section{Dispatch Latency Breakdown}

\begin{center}
\small
\begin{tabular}{@{}lrl@{}}
\toprule
\textbf{Layer} & \textbf{Median} & \\
\midrule
\texttt{cuModuleLoadData} (cold)   & 42.7 & $\mu$s \\
\texttt{cuModuleLoadData} (warm)   & 10.0 & $\mu$s \\
\texttt{cuModuleGetFunction}       &   60 & ns \\
\texttt{cuLaunchKernel} (submit)   &  1.6 & $\mu$s \\
Hot-path (launch + sync)           &  4.0 & $\mu$s \\
Binary selection logic             &    6 & ns \\
\bottomrule
\end{tabular}
\end{center}

\vspace{1pt}
{\scriptsize GTX 1650 (sm\_75), null kernel, CUBIN pre-compiled.
Cold = \texttt{execve}-isolated trials. Selection = variant table scan.}

\section{MLIR Syntax: \texttt{\#gpu.runtime\_select}}

\begin{lstlisting}
gpu.binary @kernels [
  #gpu.object<#nvptx, "sm_75.cubin">,
  #gpu.object<#amdgcn, "gfx90a.hsaco">
] { select = #gpu.runtime_select }
\end{lstlisting}

\section{OffloadBinary Metadata (5 New Keys)}

\begin{center}
\small
\begin{tabular}{@{}llp{3.2cm}@{}}
\toprule
\textbf{Key} & \textbf{Tier} & \textbf{Semantics} \\
\midrule
\texttt{min\_sm}             & MUST & Min CUDA compute cap \\
\texttt{min\_gfx}            & MUST & Min AMD GFX + family \\
\texttt{requires\_features}  & MUST & Capability tokens \\
\texttt{variant\_priority}   & MAY  & Ranking (higher = preferred) \\
\texttt{variant\_tag}        & MAY  & Label: \texttt{optimized}, \texttt{fallback} \\
\bottomrule
\end{tabular}
\end{center}

\vspace{1pt}
{\scriptsize MUST = runtime rejects if violated.
MAY = ranking only. Missing keys = no constraint (backward-compatible).}

\section{Links}

\begin{tabular}{@{}lp{3.5cm}@{}}
{[QR: GitHub]}        & Source, benchmarks, patch \\
{[QR: Discourse RFC]} & RFC discussion thread \\
\end{tabular}

\end{multicols}

\vspace*{\fill}
\hrule height 0.3pt
\vspace{2pt}
\begin{center}
\small EuroLLVM Dublin 2026 \enspace---\enspace S.\ Akash (akash@iitpatna.ac.in)
\end{center}

\end{document}
